% ---
% id: receiver/detector-imperfections/rem:noise-floor-receiver
% kind: rem
% title: "What a receiver declaration needs for a rate that rests on its model"
% summary: >-
%   Declare the loss and dark counts the device surely has, in every port;
%   then every element is positive definite and no snap revives a zero row.
%   A run that uses no enclosed reference rests, at a cutoff of two or more,
%   on its snap; an angle decides whether the reference can be exact.
% prominence: standard
% status: accepted
% uses: [prop:noise-floor, def:dark-composed, cor:exact-dark-composed, receiver/def:exact-route, receiver/thm:exact-blocks, receiver/def:complex-enclosure, foundations/rem:noise-floor, foundations/lem:rational-obstruction, foundations/def:evaluated-off-snap]
% mentions: [reductions/lem:row-relaxation, receiver/component-catalogue/def:catalogue-entries, receiver/req:exact-route-decision, foundations/req:snap-reference, foundations/req:support-safe-snap, meat-core/frank-wolfe/rem:fo-arithmetic, meat-core/frank-wolfe/req:certified-evaluation]
% assumes: []
% refs: []
% provenance: newly-derived
% lean:
% ---

The remark ends with reductions that apply a channel to the model before
comparing it with a lattice family; they are planned beyond these contracts
(\repodoc{spec/future-scope.md}).

\begin{remark}\label{rem:noise-floor-receiver}
The emitted blocks of a receiver are snapped onto a lattice before a key
rate is computed, and the dark-composed reference
(Corollary~\ref{cor:exact-dark-composed}), enclosed for any declared
parameters and exact where its radii are zero, is handed beside the snap.
Where the run's first-order bound is evaluated on that reference and its key
length is formed from a number on it, the blocks carrying it are
\tref{evaluated off its snap} (\xref{foundations/def:evaluated-off-snap}) and
the key rate rests on the declared model, not on the snap. Where the key
length is formed instead from a number certified on the snapped declaration
(\xref{meat-core/frank-wolfe/req:certified-evaluation}\,(i)), the moved
ones among those blocks rest on the snapped-measurement premise, as every
other moved block does
(\xref{foundations/req:snap-reference}\,(iii)).
This remark says what that means for a receiver declaration.

\emph{What to declare.} Model the loss and the dark counts the device has,
in every port: an efficiency $\eta_k<1$ (or any loss that leaves
$\lmax(A_\Sigma)<1$) and a dark-count probability $0<p_{\dc|k}<1$. By
Proposition~\ref{prop:noise-floor} every element of the dark-composed family
is then positive definite on every photon-number block. At a setting $t$
whose operator and prefactor in the setting family are $B^{(t)}$ and
$\gamma^{(t)}$, the least eigenvalue on block $n$ is at least
$r_S\gamma^{(t)}\big(1-\lmax(B^{(t)})\big)^{n}$ for a click pattern $S$, and
at least the sum of these over a fibre for an emitted outcome
(Proposition~\ref{prop:noise-floor}\,(i)). For the unconditioned receiver,
$B=A_\Sigma$ and $\gamma=1$, so the bound is
$r_S\big(1-\lmax(A_\Sigma)\big)^{n}$, which is at least
$r_S(1-\max_k\eta_k)^{n}$ when the loss sits in the terminal efficiency layer
(Proposition~\ref{prop:noise-floor}\,(iii)); a setting whose operator has
$\lmax(B^{(t)})\le\lmax(A_\Sigma)$ and whose prefactor $\gamma^{(t)}$ is below
one has at least $\gamma^{(t)}$ times that floor. No row of such a block is
structurally zero, so even the last resort of
\xref{foundations/req:support-safe-snap}\,(iv) gives no element of it
support that the data deny. The same floor is what makes domination by such
a model possible at a small weight
(\xref{foundations/rem:noise-floor}); no run relates a snap to its model
that way. The rates to declare are the smallest ones the
user is sure the device has. A larger rate gives a larger floor, but a model
noisier than the device is a wrong model, which nothing in the engine can
detect.

\emph{The enclosure.} An enclosed reference is evaluated at its centres,
which displaces the measured family by up to the enclosure's radius
(\xref{meat-core/frank-wolfe/rem:fo-arithmetic}). The radii of the
\tref[square-root enclosure]{square-root enclosures} and
\tref[angle enclosure]{angle enclosures}
(Definition~\ref{def:complex-enclosure}) from which the blocks are computed
are amplified on the way to the radius of a block: through the permanents of
the Fock blocks, the inclusion--exclusion sums over click patterns and the
sums over a fibre of the coarse-graining. The radius of a block bears on a
run in two places. The companion's entries are relaxed outward by it, so a
larger radius loosens them, which never raises the bound formed on the
companion (\xref{reductions/lem:row-relaxation}); and the evaluated family
lies up to the radius from the model, a perturbation that no claim bounds
in its effect on the key rate (\xref{meat-core/frank-wolfe/rem:fo-arithmetic}).
No size of the radius is a threshold for either: the enclosures are chosen
as tight as their exact computation allows, and the record shows the
largest radius of every handed block and of the companion's datum.

\emph{A run that uses no enclosed reference.} A run may declare that enclosed
references are not used (\xref{foundations/req:snap-reference}\,(i)). Then
only the blocks whose reference is exact are handed, and every other moved
block rests on the snapped-measurement premise, in the group no usable
reference, unless the run declares the strict option
(\xref{foundations/req:snap-reference}\,(iv)), which refuses. The receiver's
reference is in general exact only on the \tref{exact route}, and there with
certainty only on the blocks of at most one photon
(Theorem~\ref{thm:exact-blocks}\,(iii)): on a block of two or more photons,
an entry between occupation vectors $\vec m,\vec m'$ with
$\vec m!\,\vec m'!$ not a square is in general a Gaussian rational times an
irrational factor such as $1/\sqrt2$, and is enclosed with a positive radius.
In general, therefore, every run with a cutoff $N\ge2$ that uses no enclosed
reference rests on its snap there, and so does every such run whose
declaration is off the exact route; its record and a report on it say so.

\emph{An ideal device beyond one photon.} Without loss the no-click element
vanishes on every block $n\ge1$, and without dark counts every click pattern
vanishes on the vacuum block (Proposition~\ref{prop:noise-floor}\,(iv),
(v)); so an ideal device has rank-deficient elements on every block, and a
snap of it may need its last resort, which the record flags. A receiver that
measures in a rotated basis can have some of them rank-deficient along
irrational directions: on the one-photon block where a direction of the
basis has no Gaussian-rational representative, and on blocks of more photons
even where every direction has one. In an ideal passive BB84 receiver, the
two-photon element in which only the detector of the diagonal polarization
clicks, behind a half-wave plate at $\pi/8$, is a multiple of the projector
onto $\tfrac12\ket{2,0}+\tfrac1{\sqrt2}\ket{1,1}+\tfrac12\ket{0,2}$ in the
occupation basis of the channel's two polarization modes, and that ray
contains no Gaussian-rational vector, although the one-photon directions of
the diagonal basis are rational. By
\xref{foundations/lem:rational-obstruction}, every weight at which the model
dominates a lattice family declared in this basis is then at least the
largest eigenvalue of that element, at every lattice exponent, which is $1/4$
for a balanced basis choice; that is why a declared device weight, and not a
snap, is where domination could serve. At a cutoff $N\le1$ on the
\tref{exact route} the blocks are known exactly
(Theorem~\ref{thm:exact-blocks}\,(iii)); where they are moreover dyadic
rationals of bounded denominator, as at a balanced basis choice with dyadic
efficiencies, they are lattice values and the snap keeps them. The
obstruction belongs to domination, which compares the model with a lattice
family element by element in the occupation basis. A reduction that first
maps the $m$-photon block of the two polarization modes isometrically into
the $m$-fold tensor power of $\mathbb{C}^{2}$ carries the ray above to a
vector with rational coordinates (\xref{foundations/rem:noise-floor}); for
an ideal device on these two modes whose port Grams are Gaussian-rational it
would give a target with entries in $\mathbb{Q}(i)$ at every cutoff, at
dimension $2^{m}$ instead of $m+1$ (with $d$ input modes, $d^{m}$ instead of
$\binom{d+m-1}{m}$). The images of the elements sum to the projector onto
the symmetric subspace, not to the identity, so such a target is complete
only after an outcome is added on the complement of that subspace. No claim
of these contracts makes that reduction.

\emph{Angles.} A half-wave plate at axis angle $\theta$ has entries
$\pm\cos2\theta$ and $\pm\sin2\theta$
(\eqref{eq:hwpqwp} of Definition~\ref{def:catalogue-entries}). Every angle
is enclosed, so the angle decides nothing about whether a reference exists;
it decides whether the reference can be exact. Behind an ideal polarizing
beamsplitter the plate's port Gram on the two polarization modes has the
diagonal entries $(1\pm\cos4\theta)/2$, up to a rational factor, and the exact
route needs Gaussian-rational Grams (Definition~\ref{def:exact-route}), so it
needs $\cos4\theta$ rational. At multiples of $\pi/8$,
$\cos4\theta\in\{0,\pm1\}$; at $\theta=\pi/24$, $\cos4\theta=\sqrt3/2$ and
the reference is enclosed, which the run hands and evaluates as it does an
exact one. The \tref{exact route} reads an angle given as a rational
coefficient of $\pi$ (Requirement~\ref{req:exact-route-decision}\,(ii)); an
angle given in radians is a dyadic rational, and a non-zero one is no
rational multiple of $\pi$, so a plate declared at the binary64 value nearest
$\pi/8$ is enclosed, while one declared as $\tfrac18\pi$ can be on the
exact route.
\end{remark}
